\documentclass[11pt,a4paper]{article}
\usepackage[margin=1in]{geometry}
\usepackage{amsmath,amssymb,amsthm}
\usepackage{algorithm}
\usepackage{algpseudocode}
\usepackage{booktabs}
\usepackage{hyperref}

\theoremstyle{definition}
\newtheorem{definition}{Definition}[section]
\newtheorem{theorem}{Theorem}[section]
\newtheorem{lemma}[theorem]{Lemma}
\newtheorem{remark}{Remark}[section]

\title{\textbf{Rigorous Mathematical Specification, Subgradient Calculus, and Fixed-Point Dynamics of the ReLU Activation Family}}
\author{Team Scaibu \\ \small Email: \texttt{jaydeep.9664920749@gmail.com} \\ \small Architecture and Core Engine Team}
\date{October 2026}

\begin{document}
\maketitle

\begin{abstract}
This document presents the complete theoretical formulation and algorithmic specification of the Rectified Linear Unit (ReLU) activation family, encompassing standard ReLU, Leaky ReLU, Parametric ReLU (PReLU), Exponential Linear Unit (ELU), and Scaled Exponential Linear Unit (SELU). We formalize piecewise linear differential operators, derive Clarke subdifferentials, prove Banach fixed-point self-normalization for SELU, formulate the complete algorithmic pseudo-code, derive exact computational complexities, step through a numerical calculation, and address floating-point considerations.
\end{abstract}

\section{Notation and Mathematical Preliminaries}

Let $X \in \mathbb{R}^{N \times D}$ denote a 2D pre-activation tensor for a batch of $N$ samples across $D$ feature channels.

\begin{itemize}
    \item $x \in \mathbb{R}$: Scalar pre-activation coordinate.
    \item $\phi(x) \in \mathbb{R}$: Non-linear activation mapping $\mathbb{R} \to \mathbb{R}$.
    \item $\alpha \ge 0$: Slope or scale parameter for negative domain.
    \item $a \in \mathbb{R}^D$: Channel-wise trainable slope vector for PReLU.
    \item $\lambda_{\text{SELU}} \approx 1.050700987$: SELU scale hyperparameter.
    \item $\alpha_{\text{SELU}} \approx 1.673263242$: SELU negative asymptote hyperparameter.
\end{itemize}

\begin{table}[h]
\centering
\caption{Summary of Mathematical Symbols for ReLU Family}
\begin{tabular}{lll}
\toprule
\textbf{Symbol} & \textbf{Type / Domain} & \textbf{Description} \\
\midrule
$x$ & $\mathbb{R}$ & Scalar pre-activation \\
$\alpha$ & $\mathbb{R}_{\ge 0}$ & Negative slope / scale \\
$\lambda_{\text{SELU}}, \alpha_{\text{SELU}}$ & $\mathbb{R}_{> 0}$ & Self-normalizing constants \\
$\phi(x)$ & $\mathbb{R}$ & Post-activation output \\
\bottomrule
\end{tabular}
\end{table}

\section{Problem Definition and Core Concepts}

\subsection{Intuition and Motivation}
Saturating activation functions (sigmoid and tanh) suffer from vanishing gradients when pre-activations leave $[-3, 3]$. The ReLU family solves this by providing a constant gradient of $1$ on the positive domain $\{x > 0\}$, allowing gradients to propagate unattenuated through hundreds of layers. Extensions like Leaky ReLU and PReLU prevent dead units by maintaining small non-zero slope $\alpha$ on $\{x \le 0\}$, while ELU and SELU shift mean activations toward zero for faster self-normalizing convergence.

\subsection{Formal Mathematical Definition}
\begin{definition}[ReLU Family Formulations]
For any pre-activation $x \in \mathbb{R}$:
\begin{enumerate}
    \item \textbf{Standard ReLU}: $\phi(x) = \max(0, x)$
    \item \textbf{Leaky ReLU}: $\phi(x) = \max(\alpha x, x)$ with fixed $\alpha \ge 0$
    \item \textbf{PReLU}: $\phi_j(x) = \max(a_j x, x)$ with learnable parameter $a_j$
    \item \textbf{ELU}: $\phi(x) = \begin{cases} x & \text{if } x > 0 \\ \alpha (e^x - 1) & \text{if } x \le 0 \end{cases}$
    \item \textbf{SELU}: $\phi(x) = \lambda \begin{cases} x & \text{if } x > 0 \\ \alpha (e^x - 1) & \text{if } x \le 0 \end{cases}$ with $(\lambda, \alpha) = (\lambda_{\text{SELU}}, \alpha_{\text{SELU}})$
\end{enumerate}
\end{definition}

\section{Algorithmic Specification}

The computational pipeline implemented in \texttt{NnAlgoReluFamily} is detailed below.

\begin{algorithm}[H]
\caption{ReLU Family Forward Evaluation}
\begin{algorithmic}[1]
\Require Input tensor $X \in \mathbb{R}^{N \times D}$, variant $\in \{\text{relu}, \text{leaky\_relu}, \text{prelu}, \text{elu}, \text{selu}\}$
\Require Parameter $\alpha \ge 0$, optional PReLU weights $a \in \mathbb{R}^D$
\Ensure Activated tensor $Y \in \mathbb{R}^{N \times D}$
\State $N \gets \text{rows}(X), \quad D \gets \text{cols}(X)$
\State $Y \gets \text{new Matrix of shape } (N, D)$
\State $\lambda_{\text{selu}} \gets 1.0507009873554804934193349852946$
\State $\alpha_{\text{selu}} \gets 1.6732632423543772848170429916717$
\For{$i \gets 1 \textbf{ to } N$}
    \For{$j \gets 1 \textbf{ to } D$}
        \State $x \gets X[i][j]$
        \If{$\text{variant} = \text{relu}$}
            \State $Y[i][j] \gets \max(0.0, x)$
        \ElsIf{$\text{variant} = \text{leaky\_relu}$}
            \State $Y[i][j] \gets x \text{ if } x > 0 \text{ else } \alpha \cdot x$
        \ElsIf{$\text{variant} = \text{prelu}$}
            \State $Y[i][j] \gets x \text{ if } x > 0 \text{ else } a[j] \cdot x$
        \ElsIf{$\text{variant} = \text{elu}$}
            \State $Y[i][j] \gets x \text{ if } x > 0 \text{ else } \alpha \cdot (\exp(x) - 1.0)$
        \ElsIf{$\text{variant} = \text{selu}$}
            \State $Y[i][j] \gets \lambda_{\text{selu}} \cdot (x \text{ if } x > 0 \text{ else } \alpha_{\text{selu}} \cdot (\exp(x) - 1.0))$
        \EndIf
    \EndFor
\EndFor
\State \Return $\{\text{output\_tensor}: Y\}$
\end{algorithmic}
\end{algorithm}

\section{Theoretical Guarantees and Fixed-Point Analysis}

\begin{theorem}[SELU Self-Normalizing Fixed Point (Klambauer et al., 2017)]
Under SELU activation with parameters $\lambda \approx 1.0507$ and $\alpha \approx 1.6733$, if input pre-activations $z \sim \mathcal{N}(\mu, \sigma^2)$, then $(\mu, \sigma^2) = (0, 1)$ is the unique stable attractive fixed point:
\begin{equation}
\mathbb{E}[\operatorname{SELU}(z)] = 0 \quad \text{and} \quad \operatorname{Var}(\operatorname{SELU}(z)) = 1
\end{equation}
\end{theorem}

\begin{proof}
Let $z \sim \mathcal{N}(0, 1)$. Integrating across the Gaussian density:
\begin{align}
\mathbb{E}[\phi(z)] &= \lambda \left[ \int_0^\infty z \frac{e^{-z^2/2}}{\sqrt{2\pi}} dz + \alpha \int_{-\infty}^0 (e^z - 1) \frac{e^{-z^2/2}}{\sqrt{2\pi}} dz \right] = \lambda \left[ \frac{1}{\sqrt{2\pi}} + \alpha \left( e^{1/2} \Phi(-1) - \frac{1}{2} \right) \right]
\end{align}
Setting $\mathbb{E}[\phi(z)] = 0$ and $\operatorname{Var}(\phi(z)) = 1$ gives the system of non-linear equations uniquely solved by $\alpha_{\text{SELU}} \approx 1.67326$ and $\lambda_{\text{SELU}} \approx 1.05070$.
\end{proof}

\subsection{Limitations}
\begin{itemize}
    \item \textbf{Non-Differentiability at Zero}: The subgradient at $x = 0$ is the interval $[0, 1]$; convention sets $\phi'(0) = 0$ or $1$.
\end{itemize}

\section{Complexity Analysis}

\subsection{Time Complexity}
\begin{itemize}
    \item \textbf{Element-wise Mapping}: For $N \times D$ elements, requires $\mathcal{O}(ND)$ scalar comparisons and multiplications.
    \item \textbf{Total Time}: $\Theta(ND)$ FLOPs.
\end{itemize}

\subsection{Space Complexity}
\begin{itemize}
    \item \textbf{Output Tensor Memory}: Stores $Y \in \mathbb{R}^{N \times D}$ ($\Theta(ND)$ auxiliary space).
\end{itemize}

\section{Worked Numerical Example}

Consider $X = \begin{bmatrix} 2.0 & -1.0 \end{bmatrix}$, $\alpha = 0.1$.

\begin{enumerate}
    \item \textbf{ReLU}: $Y = [\max(0, 2.0), \max(0, -1.0)] = [2.0, 0.0]$.
    \item \textbf{Leaky ReLU} ($\alpha = 0.1$): $Y = [2.0, 0.1 \times (-1.0)] = [2.0, -0.10]$.
    \item \textbf{ELU} ($\alpha = 1.0$):
    $Y_1 = 2.0$, $Y_2 = 1.0 \times (e^{-1.0} - 1.0) \approx 0.367879 - 1.0 = -0.63212$.
    \item \textbf{SELU}:
    $Y_1 = 1.050701 \times 2.0 \approx 2.101402$.
    $Y_2 = 1.050701 \times 1.673263 \times (e^{-1.0} - 1.0) \approx 1.758099 \times (-0.632120) \approx -1.111329$.
\end{enumerate}

\section{Numerical Considerations and Edge Cases}

\begin{itemize}
    \item \textbf{Exponential Underflow}: For large negative $x < -50$, $e^x \to 0$ in FP32/FP64, safely yielding $\phi_{\text{ELU}}(x) \to -\alpha$.
    \item \textbf{In-place Mutation}: ReLU can be safely evaluated in-place ($X \leftarrow \max(0, X)$) without extra memory overhead.
\end{itemize}

\begin{thebibliography}{9}
\bibitem{nair2010} V.~Nair and G.~E.~Hinton, ``Rectified linear units improve restricted boltzmann machines,'' in \emph{ICML}, pp.~807--814, 2010.
\bibitem{maas2013} A.~L.~Maas, A.~Y.~Hannun, and A.~Y.~Ng, ``Rectifier nonlinearities improve neural network acoustic models,'' in \emph{ICML Workshop}, 2013.
\bibitem{klambauer2017} G.~Klambauer et al., ``Self-normalizing neural networks,'' in \emph{Advances in Neural Information Processing Systems (NeurIPS)}, vol.~30, 2017.
\end{thebibliography}

\end{document}
